%Cosmic rays are charged subatomic particles and nuclei originating from space and can be accelerated .
%Yet the astrophysical sources and acceleration mechanisms responsible for accelerating cosmic rays to tens of PeV and beyond are still unknown.

%These particles can be accelerated up to the energy of a baseball thrown at roughly \qty{100}{\kilo\meter\per\hour} contained within a single subatomic particle or nucleus, making them the highest energy particles in the known Universe.

%Cosmic rays are charged particles of astrophysical origin, where the highest energy cosmic rays are the most energetic particles in the known Universe.


Cosmic rays are charged particles of astrophysical origin that can be accelerated to the highest particle energies in the known Universe.
Yet, the sources and acceleration mechanisms responsible for accelerating cosmic rays to these energies, at tens of PeV and beyond, are still unknown.
Cosmic rays of these energies arriving at Earth will initiate extensive air showers, cascades of particles and radiation, complicating the determination of their origin.
Measurements of these air showers must reconstruct properties of the initial cosmic ray that hint at its source, such as the cosmic-ray energy, mass, and arrival direction.
The cosmic-ray mass in particular is correlated with two independent components of the air shower that propagate through the atmosphere to Earth's surface, the air-shower electromagnetic component and the air-shower muonic component.

This dissertation will combine independent measurements of these two shower components to study their combined effect on cosmic-ray mass sensitivity.
Air-shower measurements with radio antennas, which probe the electromagnetic component of the shower, will be combined with measurements of the muonic shower component from particle detectors.
The prospects of this hybrid technique are first explored in a simulation study which includes several additional air-shower observables.
Results of the study highlight the potential for combining radio and muon measurements to increase the cosmic-ray mass sensitivity, where these results were proven to be agnostic to the choice of hadronic interaction model.
This work then builds upon a previous analysis that predicts a muon observable from water-Cherenkov detector measurements using a neural network.
In particular, the predictive range of the network is lowered by one decade in cosmic-ray energy to near the Galactic-to-extragalactic transition region of cosmic-ray origin.
These predictions are then applied to real air-shower data and combined with a high-quality radio reconstruction of the same events to study the cosmic-ray mass sensitivity of hybrid radio-muon measurements.
Finally, a prototype station and pathfinder array of SKALA radio antennas deployed at the location of the Pierre Auger Observatory are introduced.
The pathfinder array, the Auger Radio Infill SKALA Extension (ARISE), will build upon the hybrid radio-muon technique in the coming years.


%Finally, the next steps of building upon hybrid radio-muon measurements are discussed in the context of the Auger Radio Infill SKALA Extension (ARISE)

%a neural network based predictions of 

%neural-network based muon predi

%Results of the study highlight the po were proven to be agnostic 

%based on the air-shower secondary particles and radiation

%that propagate through the atmosphere to Earth's surface.
%Measurements of these air showers must reconstruct properties of the initial cosmic ray that hint at its origin, such as its energy, mass, and arrival direction, based on the air-shower secondary particles and radiation.
%The cosmic-ray mass in particular is correlated with two independent components of the air shower that develop in the atmosphere, the air-shower electromagnetic component and the air-shower muonic component.

%These measurements must then reconstruct properties of the initial cosmic ray that hint at its origin based on the air-shower secondary particles and radiation.

%that hint at its origin

%to ultra-high-energies, wielding the energy of a 

%charged particles accelerated to ultra-high-energies in astrophysical origin, are the highest energy particl